In this Appendix we prove the \emph{neural basis expansion approximation theorem} introduced in Section~\ref{section:training_methodology}. 
We show that \ours' hierarchical interpolation
can arbitrarily approximate infinitely long horizons ($\tau \in [0,1]$ continuous horizon), as long as the interpolating functions $g$ are defined by a projections to informed multi-resolution functions, and the forecast relationships satisfy smoothness conditions. We prove the case when $g_{w,h}(\tau)=\theta_{w,h}\phi_{w,h}(\tau)=\theta_{w,h}\mathbbm{1}\{\tau \in [2^{-w}(h-1),2^{-w}h]\}$ are piecewise constants and the inputs $\ylag\in[0,1]$. The proof for linear, spline functions and $\ylag\in[a,b]$ is analogous.

\textbf{Lemma 1.} Let a function representing an infinite forecast horizon be $\mathcal{Y}: [0, 1] \rightarrow \mathbb{R}$ a square integrable function $\mathcal{L}^{2}([0,1])$. The forecast function $\mathcal{Y}$ can be arbitrarily well approximated by a linear combination of piecewise constants:
$$V_{w}=\{\phi_{w,h}(\tau) = \phi(2^{w}(\tau-h)) \;|\; w \in {\mathbb{Z}}, h\in2^{-w}\times[0,\dots,2^{w}]\} \quad $$
where $w \in \mathbb{N}$ controls the frequency/indicator's length and $h$ the time-location (knots) around which the indicator $\phi_{w,h}(\tau) = \mathbbm{1}\{\tau \in [2^{-w}(h-1),2^{-w}h]\}$ is active. That is, $\forall \epsilon >0$, there is a $w \in \mathbb{N}$ and $\hat{\mathcal{Y}}(\tau |\ylag)=\mathrm{Proj}_{V_{w}}(\mathcal{Y}(\tau|\ylag)) \in \mathrm{Span}(\phi_{w,h})$ such that

\begin{equation}
    \int_{[0,1]}|\mathcal{Y}(\tau) - \hat{\mathcal{Y}}(\tau)| d\tau =
    \int_{[0,1]} |\mathcal{Y}(\tau) - \sum_{w,h} \theta_{w,h} \phi_{w,h}(\tau) | d\tau \leq \epsilon %\quad \iff \quad \lim_{w\to\infty} |||| = 0
\end{equation}

\begin{proof}
This classical proof can be traced back to Haar's work (1910). The indicator functions $V_{w} = \{\phi_{w,h}(\tau)\}$ are also referred in literature as Haar scaling functions or father wavelets. Details provided in \cite{boggess2015first}.Let the number of coefficients for the $\epsilon$-approximation $\hat{\mathcal{Y}}(\tau |\ylag)$ be denoted as $N_{\epsilon} = \sum^{w}_{i=0} 2^i $.

\end{proof}


\textbf{Lemma 2.}
Let a forecast mapping $\mathcal{Y}(\cdot\;|\;\ylag): [0,1]^{L} \to \mathcal{L}^{2}([0,1])$ be $\epsilon$-approximated by $\hat{\mathcal{Y}}(\tau |\ylag)=\mathrm{Proj}_{V_{w}}(\mathcal{Y}(\tau|\ylag))$, the projection to multi-resolution piecewise constants. If the relationship between $\ylag \in [0,1]^{L}$ and $\theta_{w,h}$ varies smoothly, for instance $\theta_{w,h}:[0,1]^{L}\to \mathbb{R}$ is a K-Lipschitz function then for all $\epsilon > 0$ there exists a three-layer neural network $\hat{\theta}_{w,h}: [0, 1]^L \rightarrow \mathbb{R}$ with $O\left(L\right(\frac{K}{\varepsilon}\left)^L\right)$ neurons and $\mathrm{ReLU}$ activations such that 
\begin{equation}
\int_{[0, 1]^L} |\theta_{w,h}(\ylag) - \hat{\theta}_{w,h}(\ylag)| d\ylag \leq \epsilon
\end{equation}

\begin{proof}
This lemma is a special case of the neural universal approximation theorem that states the approximation capacity of neural networks of arbitrary width \citep{hornik1991approximation}. The theorem has refined versions where the width can be decreased under more restrictive conditions for the approximated function \citep{barron1993approximation, boris2017approximation}.
\end{proof}

% \vspace{5mm}
\newpage
\textbf{Theorem 1.} Let a forecast mapping be

$\mathcal{Y}(\cdot \;|\; \ylag): [0, 1]^L \rightarrow \mathcal{F}$, where the forecast functions $\mathcal{F}=\{\mathcal{Y}(\tau): [0,1] \to \mathbb{R}\}=\mathcal{L}^{2}([0,1])$ representing a continuous horizon, are square integrable. 

If the multi-resolution functions $V_{w}$ can arbitrarily approximate $\mathcal{L}^{2}([0,1])$. And the projection $\mathrm{Proj}_{V_{w}}(\mathcal{Y}(\tau))$ varies smoothly on $\ylag$. Then the forecast mapping $\mathcal{Y}(\cdot \;|\; \ylag)$ can be arbitrarily approximated by a neural network learning a finite number of  multi-resolution coefficients $\hat{\theta}_{w,h}$. %$\overline{\cup_{w \in \mathbb{Z}} V_{w}} = \mathcal{L}^{2}([0,1])$

That is $\forall \epsilon>0$, 
\begin{align}
% \nonumber
\begin{split}
    \int |\mathcal{Y}(\tau \;|\; \ylag) -  \tilde{\mathcal{Y}}(\tau \;|\; \ylag) | d\tau \qquad\qquad\qquad\qquad\qquad \\
          = \int |\mathcal{Y}(\tau \;|\; \ylag) 
          -\sum_{w,h} \hat{\theta}_{w,h}(\ylag) \phi_{w,h}(\tau)| d\tau \leq \epsilon
\end{split}
\end{align}

% \begin{equation}
%     \int |\mathcal{Y}(\tau \;|\; \ylag) - \tilde{\mathcal{Y}}(\tau \;|\; \ylag) | d\tau = \int |\mathcal{Y}(\tau \;|\; \ylag) -\sum_{w,h} \hat{\theta}_{w,h}(\ylag) \phi_{w,h}(\tau)| d\tau \leq \epsilon
% \end{equation}

% \clearpage
\begin{proof} For simplicity of the proof, we will omit the conditional lags $\ylag$. Using both the neural approximation $\tilde{\mathcal{Y}}$ from Lemma 2, and Haar's approximation $\hat{\mathcal{Y}}$ from Lemma 1, 
\begin{align*}
    \int |\mathcal{Y}(\tau) -\tilde{\mathcal{Y}}(\tau)| d\tau 
    &= 
    \int |(\mathcal{Y}(\tau) - \hat{\mathcal{Y}}(\tau))+(\hat{\mathcal{Y}}(\tau)-\tilde{\mathcal{Y}}(\tau))| d\tau 
\end{align*}

By the triangular inequality:
\begin{align*}
\begin{split}
    \int |\mathcal{Y}(\tau) -\tilde{\mathcal{Y}}(\tau)| d\tau & \leq \int |\mathcal{Y}(\tau) - \hat{\mathcal{Y}}(\tau)| \\
    & +
    |\sum_{w,h} \theta_{w,h} \phi_{w,h}(\tau) - \sum_{w,h} \hat{\theta}_{w,h} \phi_{w,h}(\tau)| d \tau %\\
\end{split}
\end{align*}

By a special case of Fubini's theorem
\begin{align*}
    \int |\mathcal{Y}(\tau) -\tilde{\mathcal{Y}}(\tau)| d\tau \leq 
    \qquad\qquad\qquad\qquad\qquad\qquad\qquad\quad \\
    \int 
    |\mathcal{Y}(\tau) - \sum_{w,h} \hat{\mathcal{Y}}(\tau)| d \tau 
    +
    \sum_{w,h} \int_{\tau} 
    |(\theta_{w,h} - \hat{\theta}_{w,h}) \phi_{w,h}(\tau)| d \tau %\\
\end{align*}    

Using positivity and bounds of the indicator functions
\begin{align*}
    \int |\mathcal{Y}(\tau) -\tilde{\mathcal{Y}}(\tau)| d\tau 
    \leq \qquad\qquad\qquad\qquad\qquad\qquad\qquad\quad\\
    \int_{\tau} 
    |\mathcal{Y}(\tau) - \sum_{w,h} \hat{\mathcal{Y}}(\tau)| d \tau 
    +
    \sum_{w,h} |\theta_{w,h} - \hat{\theta}_{w,h}| \int_{\tau} \phi_{w,h}(\tau) d \tau \\
    < \int_{\tau} 
    |\mathcal{Y}(\tau) - \sum_{w,h} \hat{\mathcal{Y}}(\tau)| d \tau
    +
    \sum_{w,h} |\theta_{w,h} - \hat{\theta}_{w,h}| 
\end{align*} 

To conclude we use the both arbitrary approximations from the Haar projection and the approximation to the finite multi-resolution coefficients
\begin{align*}
    \int |\mathcal{Y}(\tau) -\tilde{\mathcal{Y}}(\tau)| d\tau    
    \leq \qquad\qquad\qquad\qquad\qquad\qquad\qquad\quad\\
    \int |\mathcal{Y}(\tau) - \hat{\mathcal{Y}}(\tau)| d \tau 
    +
    \sum_{w,h} |\theta_{w,h} - \hat{\theta}_{w,h}| %\\
    \leq \epsilon_{1} + N_{\epsilon_{1}} \epsilon_{2} \; \leq \epsilon
\end{align*}
\end{proof}

% \newpage

% \EDIT{Don't know how to use this}

% Let there be a collection of function subspaces $V_{w}$ that satisfy the following properties:

% \begin{enumerate}
%     \item Succesive approximation subspaces: $\quad \dots V_{2} \subset V_{1} \subset V_{0} \subset V_{-1} \subset V_{-2} \dots$
%     \item Anti-redundancy and completeness: $\quad \cap_{w \in \mathbb{Z}} V_{w} = \{0\}$ and $\; \overline{\cup_{w \in \mathbb{Z}} V_{w}} = \mathcal{L}^{2}[\mathbb{R}]$
%     \item Scale self-similarity: $\hat{f} \in V_{w} \implies \hat{f}(2^{w}\cdot) \in V_{0} \quad \forall w\in\mathbb{Z}$
%     \item Time self-similarity: $\hat{f} \in V_{0} \implies \hat{f}(\cdot - h) \in V_{0} \quad \forall h\in\mathbb{Z}$
%     \item Regularity: $\exists \phi \in V_{0}$, such that its integer shifts $\{\phi_{0,h} \;|\; h \in \mathbb{Z}\}$ are an orthonormal basis of $V_{0}$
% \end{enumerate}

% Then there exists an orthonormal wavelet basis $\{\psi_{w,h}(\tau) = 2^{-w/2}\psi(2^{-w}\tau-h) \;|\; w,h\in\mathbb{Z} \}$ such that $\forall f \in \mathcal{L}^{2}(\mathbb{R})$
% \begin{align}
%   \mathrm{Proj}_{V_{w-1}} f = \mathrm{Proj}_{V_{w}} f + \sum_{h} \langle f, \psi_{w,h} \rangle \psi_{w,h} \\
%   V_{w} = V_{w} \oplus V^{\bot}_{w} \qquad \mathcal{L}^{2}(\mathbb{R}) = \overline{ \bigoplus_{w \in \mathbb{Z}} V^{\bot}_{z} }
% \end{align}